% This is a template for students in MATH 3000 at FSU to use for the final draft of their final exam.

% All of this stuff with '%' in front is a comment and ignored by the compiler.
%
% The lines before the "\begin{document}" line is called the preamble.
% This is where you load particular packages you need.
% Until you are more experienced, or the program says you are missing packages, it is safe to ignore most of the preamble.
%
%----------------------------------

\documentclass[12pt]{article}
\usepackage[margin=1in]{geometry}% Change the margins here if you wish.
\setlength{\parindent}{0pt} % This is the set the indent length for new paragraphs, change if you want.
\setlength{\parskip}{5pt} % This sets the distance between paragraphs, which will be used anytime you have a blank line in your LaTeX code.
\pagenumbering{gobble}% This means the page will not be numbered. You can comment it out if you like page numbers.

%------------------------------------

% These packages allow the most of the common "mathly things"
\usepackage{amsmath,amsthm,amssymb}

% This package allows you to add images.
\usepackage{graphicx}
\usepackage{float}

% These are theorem environments.  This should cover everything you need, and you should be able to tell what environment goes with what type of result, but please let me know if I've missed anything.
\newtheorem{euclidtheorem}{Proposition}[section]
\newtheorem{classtheorem}{Theorem}
\newtheorem{theorem}{Theorem}[section]
\newtheorem*{theorem*}{Theorem}


%These help to format the names of the results the way we are in class and in notes.
\renewcommand*{\theeuclidtheorem}{\Roman{section}.\arabic{theorem}}
\renewcommand*{\thetheorem}{\arabic{section}.\arabic{theorem}}
\renewcommand*{\theclasstheorem}{\Alph{theorem}}



% Should you need any additional packages, you can load them here. If you've looked up something (like on DeTeXify), it should specify if you need a special package.  Just copy and paste what is below, and put the package name in the { }.  
\usepackage{wasysym} %this lets me make smiley faces :-)

\title{Geometry Final}

% You are the author, put your name here.
\author{Sarah Wright}

% You can change the date to be something other than the current date if you want.
\date{\today}

\begin{document}

\maketitle


\subsection*{Introduction}
If you did an oral portion of this final exam, include a brief summary here.  Be sure to include if you successfully completed any of the problems.  If you did complete portions, you can delete those sections below.  

Below should be a formal write-up of final draft quality; follow the same Writing Style Guide as we use for papers.  You only need to complete this for problems that you did not complete during your oral exam.

\subsection*{Regular Octagon Construction}

Include three things here:
\begin{itemize}
\item  The steps for your construction of a regular octagon and pictures of (at least most of) those steps.
\item  A proof that what you have constructed is a regular octagon.
\item  A brief summary of how general your construction is.
\end{itemize}
These do not need to be organized in this exact way; you can mush some things together.  As long as I can find all the parts, you're good to go.

\subsection*{Task Involving ``Content"}
  
This section will look different depending on which task from Chapter 11 you choose to complete, but in general you should follow the guidelines for regular papers.  Include a brief intro, an appropriate statement of the result or task, and the supporting evidence.


\subsection*{Paper Folding}

For this portion, you should be writing for an audience that is unfamiliar with this question.  Explain the folding required 
% I've copied my explanation below in case that is helpful, but definitely don't just use mine.
and any special cases you noticed where the folding directions should be different.  Some results are difficult to state formally.  Include those types of things that you noticed in this narrative portion.

In a second part, include formal statements of any conjectures you have.  You should aim to move as many statements from the informal section above to this part as you can.

In the third section, include the proof of one of your conjectures.  Again, the more you have, the better, but definitely prove something.

In a final section, include any related ideas you could explore.  You don't have to actually do any of this nor make any formal or informal conjectures.  These statements could very likely start with the phrase, ``I wonder what would happen if \dots''

Start with a square piece of paper.  Make a single fold in the paper.  Any fold you wish.
\begin{center}\includegraphics[width = .35\textwidth]{square} \includegraphics[width = .35\textwidth]{motherline}\end{center}
\noindent This splits the four edges of the square into six segments.  Fold each of these pieces to meet the main blue fold:
\begin{center}\includegraphics[width = .35\textwidth]{firstfold}\includegraphics[width = .35\textwidth]{secondfold}\includegraphics[width = .35\textwidth]{allfolds}\end{center}


\end{document}